\src{179}
\result{THEOREM}{1} If $Z^{n-p}$ and $Z'^{n-p'}$ are nonsingular cycles that
intersect transversely, then
\begin{equation}\tag{4.4}
  P_X(Z.Z')=P_X(Z).P_X(Z'),
\end{equation}
where the product on the right is the cup product
\[
  H^{p}(X,\Omega_X^{p})\times H^{p'}(X,\Omega_X^{p'})
  \longrightarrow H^{p+p'}(X,\Omega_X^{p+p'}).
\]
(We assume that $X$ is isomorphic to a locally closed subset of a projective
space.)

This last hypothesis is used only to conclude that every coherent algebraic
sheaf on $X$ is a quotient of a locally free coherent algebraic sheaf (SERRE),
and hence admits a left resolution by locally free sheaves. To prove Theorem 1,
we may suppose that $Z$ and $Z'$ are irreducible nonsingular subvarieties $Y$
and $Y'$ meeting transversely. Let $L_{*}$ be a left resolution of $\Osh_Y$ by
locally free sheaves. By Proposition 3, the diagram of homomorphisms (4.3) then
identifies with
\[
R^{p}\Gamma(\Omega_X^{p})
\xleftarrow{\alpha}
\mathbb R^{p}\Gamma\!\left(\Homsh_{\Osh_X}(L_{*},\Omega_X^{p})\right)
\xrightarrow[\sim]{\beta}
\Gamma\!\left(H^{p}\bigl(\Homsh_{\Osh_X}(L_{*},\Omega_X^{p})\bigr)\right).
\]
Here $\Gamma$ is the “group of sections” functor on the category of sheaves of
abelian groups on $X$, $\mathbb R^{p}\Gamma$ denotes its hypercohomology in
dimension $p$, and $R^{p}\Gamma$ its $p$-th derived functor. For simplicity, we
may suppose that $L_0=\Osh_X$ and that the augmentation
$L_0\longrightarrow\Osh_Y$ is the natural homomorphism (which is permissible).
Then $\alpha$ is induced by the homomorphism of complexes
$\Osh_X\longrightarrow L_{*}$ (with $\Osh_X$ regarded as a complex concentrated
in degree $0$), taking account of the equality
$\mathbb R^{p}\Gamma(K)=R^{p}\Gamma(K_0)$ when $K$ is a complex of sheaves
concentrated in degree $0$. The homomorphism $\beta$ is the familiar edge
homomorphism. Consider the analogous diagram for a locally free resolution
$L'_{*}$ of $\Osh_{Y'}$, and then the following commutative diagram of pairings:
\begin{equation}\tag{4.5}
\resizebox{\textwidth}{!}{$
\begin{array}{ccc}
R^{p}\Gamma(\Omega_X^{p})
&\longleftarrow
\mathbb R^{p}\Gamma\!\left(\Homsh_{\Osh_X}(L_{*},\Omega_X^{p})\right)
&\xrightarrow{\sim}
\Gamma H^{p}\!\left(\Homsh_{\Osh_X}(L_{*},\Omega_X^{p})\right)
\\[3pt]
\times&\times&\times\\[3pt]
R^{p'}\Gamma(\Omega_X^{p'})
&\longleftarrow
\mathbb R^{p'}\Gamma\!\left(\Homsh_{\Osh_X}(L'_{*},\Omega_X^{p'})\right)
&\xrightarrow{\sim}
\Gamma H^{p'}\!\left(\Homsh_{\Osh_X}(L'_{*},\Omega_X^{p'})\right)
\\[3pt]
\downarrow&\downarrow&\downarrow\\[3pt]
R^{p+p'}\Gamma(\Omega_X^{p+p'})
&\longleftarrow
\mathbb R^{p+p'}\Gamma\!\left(
\Homsh_{\Osh_X}(L_{*}\otimes L'_{*},\Omega_X^{p+p'})\right)
&\xrightarrow{\sim}
\Gamma H^{p+p'}\!\left(
\Homsh_{\Osh_X}(L_{*}\otimes L'_{*},\Omega_X^{p+p'})\right).
\end{array}
$}
\end{equation}

